\section{Introduction}

The gonality of a projective variety $X$ is defined in this paper as the minimal gonality of a smooth projective irreducible curve $C$ admitting a nonconstant morphism $j:C\rightarrow X$. In the case of an abelian variety, the gonality is the same as the covering gonality studied in~\cite{lazarsfeld}. One of the main results of this paper answers affirmatively a question asked in~\cite{lazarsfeld} concerning the gonality of a very general abelian variety $A$, namely, whether it grows to infinity with $g=\dim A$.

\begin{theo} \label{theogon}
Let $A$ be a very general abelian variety of dimension $g$. If $g\geq 2^{k-2}(2k-1)+(2^{k-2}-1)(k-2)$, the gonality of $A$ is at least $k+1$.
\end{theo}

In other words, a very general abelian variety of dimension $\geq 2^{k-2}(2k-1)+\linebreak
(2^{k-2}-1)(k-2)$ does not contain a curve of gonality $\leq k$. This theorem is presumably not optimal. What seems reasonable is the following bound.

\begin{conj}\label{conj}
Let $A$ be a very general abelian variety of dimension $g$.
If $g\geq 2k-1$, the gonality of $A$ is at least $k+1$.
\end{conj}

We will discuss in Section~\ref{secdiscussion} a strategy towards proving this statement and some evidence for it. Theorem~\ref{theogon} generalizes the following result by Pirola~\cite{pirola}.

\begin{theo}[Pirola]\label{theopirola}
A very general abelian variety of dimension at least $3$ contains no hyperelliptic curve.
\end{theo}

We will in fact use in the proof of Theorem~\ref{theogon} (and also Theorems~\ref{coroorbit}, \ref{theochowring} below) some of the arguments in~\cite{pirola} that we generalize in Section~\ref{secpirola}. Theorem~\ref{theogon} will be obtained as a consequence of the study of $0$-cycles modulo rational equivalence on abelian varieties. This generalized setting already appears in the paper~\cite{alzatipirola} where some improvements of Theorem~\ref{theopirola} (for example on the nonexistence of trigonal curves on very general abelian varieties of dimension $\geq 4$) were obtained.

In this article, the Chow groups with $\mathbb{Q}$-coefficients of a variety $X$ are denoted by $\CH(X)$. Rational equivalence of $0$-cycles is not very well understood, despite Mumford's theorem~\cite{mumford}. The most striking phenomenon is the existence of surfaces (e.g. of Godeaux type, see~\cite{voisingodeaux}) which are of general type but have trivial $\CH_0$-group. In the papers~\cite{voisink3}, \cite{voisinisot}, we emphasized nevertheless the geometric importance, particularly in the case of $K3$ surfaces and hyper-K\"ahler manifolds, of the study of orbits of $0$-cycles $Z$ of $X$ under rational equivalence, namely
\[
| Z|=\{Z'\in X^{(k)},\,\,Z'\,\,\text{rationally equivalent to } Z \text{ in }X\}.
\]
Here $X^{(k)}$ is the $k$-th symmetric product of $X$, or equivalently the set of effective $0$-cycles of $X$ of degree $k$. This orbit is an analogue for higher dimensional varieties $X$ of the linear system $|Z|$, for a divisor $Z$ on a smooth curve.

It is a general fact that these orbits are countable unions of closed algebraic subsets in the symmetric product of the considered variety, so that their dimension is well defined. Below we denote by $\{x\}$ the $0$-cycle of a point $x\in A$ and $0_A$ will be the origin of $A$. The following results concerning orbits $|Z|\subset A^{(k)}$ for rational equivalence, and in particular the orbit $|k\{0_A\}|$, can be regarded as a Chow-theoretic version of Theorem~\ref{theogon}.

\begin{theoc}\ \label{coroorbit}
\begin{enumerate}%\romanenumi
\item \label{1.4.1} For any abelian variety $A$ and integer $k\geq 1$, any orbit $|Z|\subset A^{(k)}$ has dimension $\leq k-1$.

\item \label{1.4.2} If $g\geq 2^{k-1}(2k-1)+(2^{k-1}-1)(k-2)$, a very general abelian variety $A$ of dimension $g$ has no positive-dimensional orbit $|Z|$, with $\deg Z\leq k$.

\item \label{1.4.3} If $k\geq 2$ and $A$ is very general of dimension $g\geq 2^{k-2}(2k-1)+(2^{k-2}-1)(k-2)$, $A$ has no positive-dimensional orbit of the form $|Z'+2\{0_A\}|$, with $Z'$ effective and $\deg Z'\leq k-2$.

\item \label{1.4.4} If $A$ is a very general abelian variety of dimension $g\geq 2k-1$, the orbit $|k\{0_A\}|$ is countable.
\end{enumerate}
\end{theoc}

In fact, Theorem~\ref{coroorbit}$\MK$\eqref{1.4.3} implies Theorem~\ref{theogon}, because for a $k$-gonal irreducible curve $C$, with nonconstant morphism $j: {C}\rightarrow A$, any degree-$k$ divisor $D\in \Pic({C})$ with $h^0({C},D)\geq 2$ provides a positive-dimensional orbit $\{j_*D'\}_{D'\in |D|}$ in $A^{(k)}$ as in~\eqref{1.4.3}. Indeed, $C$ is not rational, so $k\geq2$, and we can assume that one Weierstrass point $c\in {C}$ of $|D|$, that is, a point $c$ such that $h^0({C},D(-2c))\not=0$, is mapped to $0_A$ by $j$, and this provides a positive-dimensional orbit of the form $|Z'+2\{0_A\}|$, with $Z'$ effective and $\deg Z'= k-2$.

Item~\eqref{1.4.1} of Theorem~\ref{coroorbit} will be proved in Section~\ref{secdimorbit} (cf. Theorem~\ref{ledimorbit}). The estimates in Theorems~\ref{theogon} and~\ref{coroorbit}$\MK$\eqref{1.4.2} can probably be strongly improved. Estimate~\eqref{1.4.1} in Theorem~\ref{coroorbit} cannot be improved. To start with, it is optimal for $g=1$ because for any degree-$k$ divisor $D$ on an elliptic curve $E$, we have $|D|=\mathbb{P}^{k-1}\subset E^{(k)}$. This immediately implies that the statement is optimal for any $g$ because for abelian varieties $A=E\times B$ admitting an elliptic factor, we have $E^{(k)}\subset A^{(k)}$.

In the case $g=2$, we observe that orbits $|Z|\subset A^{(k)}$ are contained in the generalized Kummer variety $K_{k-1}(A)$ constructed by Beauville~\cite{Beau}. (More precisely, this is true for the open set of $|Z|$ parameterizing cycles where all points appear with multiplicity $1$ but this is a minor point, cf.~\cite{voisink3} for a discussion of cycles with multiplicities.) This variety is of dimension~$2k-2$ and has an everywhere nondegenerate holomorphic $2$-form for which any orbit $|Z|$ is totally isotropic, which implies the estimate~\eqref{1.4.1} in the case $g=2$. Furthermore they are also orbits for rational equivalence in $K_{k-1}(A)$, as proved in~\cite{marian}, hence they are in fact constant cycle subvarieties in $K_{k-1}(A)$ in the sense of Huybrechts~\cite{huy}.

The question whether Lagrangian (that is, of maximal dimension) constant cycle subvarieties exist in hyper-K\"ahler manifolds is posed in~\cite{voisinisot}. For a general abelian variety $A$, choosing a smooth curve $C\subset A$ of genus $g'$, we have $C^{(k)}\subset A^{(k)}$ for any $k$ and $C^{(k)}$ contains linear systems $\mathbb{P}^{k-g'}$, for $k\geq g'$. So when $k$ tends to infinity, the estimate~\eqref{1.4.1} has optimal growth in $k$.

Theorem~\ref{coroorbit}$\MK$\eqref{1.4.4}, which will be proved in Section~\ref{secproofcoriv}, has the following immediate consequence (which is a much better estimate than the one given in Theorem~\ref{theogon}).

\begin{coro}
If $A$ is a very general abelian variety of dimension $g\geq 2k-1$ and $C\rightarrow A$ is any nonconstant morphism from a smooth projective irreducible curve ${C}$, one has $h^{0}({C},\mathcal{O}_{{C}}(kc))=1$ for any point $c\in {C}$.
\end{coro}

This corollary could be regarded as the right generalization of Theorem~\ref{theopirola}.

\begin{rema}
Pirola proves in~\cite{pirola2} that for a very general abelian variety $A$ of dimension $g\geq4$, any curve $C\subset A$ has genus $\geq \frac{g(g-1)}{2}+1$. This suggests that Theorem~\ref{coroorbit}$\MK$\eqref{1.4.4} is not optimal and that an inequality $g\geq O(\sqrt{k})$ should already imply the countability of $|k\{0_A\}|$.
\end{rema}

We will give two proofs of Theorem~\ref{coroorbit}$\MK$\eqref{1.4.4}. One of them will use Theorem~\ref{theochowring} and Proposition~\ref{theorelation}, which are statements of independent interest concerning the Chow ring (as opposed to the Chow groups) of an abelian variety $A$, which we now describe. Here the product in the Chow ring is the intersection product but one can also consider the ring structure given by the Pontryagin product $*$ defined by
\[
z*z'=\mu_*(z\times z')
\]
where $\mu:A\times A\rightarrow A$ is the sum map and $z\times z'=pr_1^*z\cdot pr_2^*z'$ for $z,\,z'\in \CH(A)$. The two rings are related via the Fourier transform, see~\cite{beauvillefourier}. Define
\begin{equation}\label{eqAk}
A_k\subset A
\end{equation}
to be the set of points $x\in A$ such that $(\{x\}-\{0_A\})^{*k}=0$ in $\CH_0(A)$. We can also define $\widehat{A}_k\subset \widehat{A}$ to be the set of $D\in \Pic^0(A)=:\widehat{A}$ such that $D^k=0$ in $\CH^k(A)$. These two sets are related as follows (see Section~\ref{secproofcoriv} for a proof): choose a polarization $\theta$ on $A$, that is, an ample divisor. The polarization gives an isogeny of abelian varieties
\begin{align*}
A&\rightarrow \widehat{A},\\
x&\mapsto D_x:=\theta_x-\theta.
\end{align*}

\begin{lemm}\label{lefourier}
One has $D_x^k=0$ in $\CH^k(A)$ if and only if $(\{x\}-\{0_A\})^{*k}=0$ in $\CH_0(A)$.
\end{lemm}

Our result concerning the sets $A_k$ and $\widehat{A}_k$ is the following.

\begin{theo}\label{theochowring}
Let $A$ be an abelian variety of dimension $g$. Then, for any positive integer $k$, one has
\begin{enumerate}
\item \label{1.8.1} $\dim A_k\leq k-1$ and $\dim \widehat{A}_k\leq k-1$.
\item \label{1.8.2} If $A$ is very general and $g\geq 2k-1$, the sets $\widehat{A}_k$ and $A_k$ are countable.
\end{enumerate}
\end{theo}

Note that in~\eqref{1.8.1} and~\eqref{1.8.2}, the two stated properties are equivalent by Lemma~\ref{lefourier}, since, if $A$ is very general, so is $\widehat{A}$.

The fact that Theorem~\ref{theochowring} implies Theorem~\ref{coroorbit}$\MK$\eqref{1.4.4}, uses the following intriguing result that does not seem to be written anywhere, although some related results are available, in particular in~\cite{colombo}, \cite{herbaut}, \cite{voisindiag}.

\begin{prop}\label{theorelation}
Let $A$ be an abelian variety and let $x_1,\dots,\,x_k$ be $k$ points of $A$ such that $\sum_{i=1}^k\{x_i\}-k\{0_A\}=0$ in $\CH_0(A)$. Then for any $i\in\{1,\dots,\,k\}$
\begin{equation}\label{eqrod}
(\{x_i\}-\{0\})^{*k}=0 \text{ in }\CH_0(A).
\end{equation}
In other words, $x_i\in A_k$.
\end{prop}

For the proof of Theorem~\ref{theochowring}, we will show how the dimension estimate provided by~\eqref{1.8.1} implies the nonexistence theorem stated in~\eqref{1.8.2}. This is obtained by establishing and applying Theorem~\ref{thegenpirola}, which we will state in Section~\ref{secpirola}. This theorem, which is obtained by a direct generalization of Pirola's arguments in~\cite{pirola}, says that ``naturally defined subsets'' of abelian varieties (see Definition~\ref{definat}), assuming they are proper subsets for very general abelian varieties of a given dimension $g$, are at most countable for very general abelian varieties of dimension $\geq 2g-1$.


\section{Naturally defined subsets of abelian varieties}\label{secpirola}

The proof of Theorem~\ref{theopirola} by Pirola has two steps. First of all, Pirola shows that hyperelliptic curves in an abelian variety $A$, one of whose Weierstrass points coincides with $0_A$, are rigid. Secondly, he deduces from this rigidity statement the nonexistence of any hyperelliptic curve in a very general abelian variety of dimension $\geq 3$ by an argument of specialization to abelian varieties isogenous to a product $B\times E$, that we now extend to cover more situations.

\begin{defi}\label{definat}
We will say that the data of attaching a subset $\Sigma_A\subset A$ to any abelian variety $A$ provide naturally defined subsets if they satisfy
the following conditions:
\begin{enumerate}\setcounter{enumi}{-1}
\item \label{2.1.0} $\Sigma_A\subset A$ is a countable union of closed algebraic subsets of $A$.
\item \label{2.1.1} For any morphism $f:A\rightarrow B$ of abelian varieties, $f(\Sigma_A)\subset \Sigma_B$.
\item \label{2.1.2} For any family $\mathcal{A}\rightarrow S$, there is a countable union of closed algebraic subsets $\Sigma_{\mathcal{A}}\subset \mathcal{A}$ whose set-theoretic fibers over $S$ satisfy $\Sigma_{\mathcal{A},b}=\Sigma_{\mathcal{A}_b}$, for any $b\in S$.
\end{enumerate}
\end{defi}

\begin{rema}
By a morphism of abelian varieties $A,\,B$, we mean a group morphism, that is, mapping $0_A$ to $0_B$.
\end{rema}

Recall that the dimension of a countable union of closed algebraic subsets is defined as the supremum of the dimensions of its components (which are well defined since we are over the uncountable field $ \mathbb{C}$).

\begin{theo} \label{thegenpirola}
Let $\Sigma_A\subset A$ be naturally defined.
\begin{enumerate}
\item \label{2.3.1} Assume that for any abelian variety $A$ of dimension $g_0$, one has $\Sigma_A\not=A$. Then for a very general abelian variety $A$ of dimension $\geq 2g_0-1$, $\Sigma_A$ is at most countable.

\item \label{2.3.2} Assume that $\dim \Sigma_A\leq k$ for any $A$. Then for a very general abelian variety $A$ of dimension $\geq 2k+1$, the set $\Sigma_A$ is at most countable.

\item \label{2.3.3} Assume that $\dim \Sigma_A\leq k-1$ for a very general abelian variety $A$ of dimension $g_0\geq k$. Then for a very general abelian variety $A$ of dimension $\geq g_0+k-1$, the set $\Sigma_A$ is at most countable.
\end{enumerate}
\end{theo}

Statement~\eqref{2.3.2} is a particular case of~\eqref{2.3.1} where we take $g_0=k+1$. Both~\eqref{2.3.1} and~\eqref{2.3.3} will follow from the following result:

\begin{propc}\ \label{propourtheogenpirola}
\begin{enumerate}
\renewcommand{\theenumi}{a}
\item \label{2.4.a}
If for a very general abelian variety $B$ of dimension $g> k$, one has
\[
\dim \Sigma_B\leq k
\]
then, for a very general abelian variety $A$ of dimension $g+1$, one has
\[
\dim \Sigma_A\leq k.
\]

\renewcommand{\theenumi}{a$'$}
\item \label{2.4.a'}
In particular, if the set $\Sigma_B$ is countable for a very general abelian variety $B$ of dimension $g>0$, then for a very general abelian variety $A$ of dimension $\geq g$, the set $\Sigma_A$ is countable.

\renewcommand{\theenumi}{b}
\item \label{2.4.b}
In the situation of~\eqref{2.4.a}, if furthermore
\[
0<\dim \Sigma_B\leq k
\]
then, for a very general abelian variety $A$ of dimension $g+1$, one has
\[
\dim \Sigma_A\leq k-1.
\]
\end{enumerate}
\end{propc}

Indeed, applying Proposition~\ref{propourtheogenpirola}, we conclude in case~\eqref{2.3.1} that the dimension of $\Sigma_A$ for $A$ very general is decreasing with $g\geq g_0$ and strictly decreasing as long as it is not equal to $0$, and by assumption it is not greater than $g_0-1$ for $g=g_0$. Hence this dimension must be at most $0$ (that is, $\Sigma_A$ is at most countable) for some $g\leq 2g_0-1$. By Proposition~\ref{propourtheogenpirola}$\MK$\eqref{2.4.a'}, we then conclude that $\Sigma_A$ is countable for $A$ very general for any $g\geq 2g_0-1$.

In case~\eqref{2.3.3} of the theorem, the argument is the same except that we start with dimension $g_0=k+1$ and we conclude similarly that the dimension of $\Sigma_A$ for very general $A$ is strictly decreasing with $g\geq g_0$ as long as it is not equal to $0$. Furthermore, for $g=g_0$, this dimension is at most $k-1$. Hence the dimension of $\Sigma_A$ for very general $A$ must be $0$ for some $g\leq g_0+k-1$ and thus, by Proposition~\ref{propourtheogenpirola}$\MK$\eqref{2.4.a'}, $\Sigma_A$ is countable for any $g\geq g_0+k-1$ and $A$ very general. This proves Theorem~\ref{thegenpirola} assuming Proposition~\ref{propourtheogenpirola} that we now prove along the same lines as in~\cite{pirola}.

\begin{proof}[Proof of Proposition~\ref{propourtheogenpirola}]
Assume that $\dim \Sigma_A= k'$ for a very general abelian variety $A$ of dimension $g+1$. From the definition of naturally defined subsets, and by standard arguments involving the properness and countability properties of relative Chow varieties, there exist, for each universal family $\mathcal{A}\rightarrow S$ of polarized abelian varieties with given polarization type $\theta$, a generically finite dominant base-change morphism $S'\rightarrow S$, and, denoting by $\mathcal{A}_{S'}\rightarrow S'$ the base-changed family, a closed algebraic subset
\[
\Sigma'_{\mathcal{A}}\subset \Sigma_{\mathcal{A}_{S'}}\subset\mathcal{A}_{S'},
\]
such that the morphism $\Sigma'_{\mathcal{A}}\rightarrow S'$ is flat with irreducible fibers of dimension $k'$. In other words, we choose one $k'$-dimensional component of $\Sigma_A$ for each $A$, and we can do this in families, maybe after passing to a generically finite cover of a Zariski open set of the base.

The main observation is now the fact that there is a dense countable union of algebraic subsets $S'_\lambda\subset S'$ along which the fiber $\mathcal{A}_b$ is isogenous to a product $B_\lambda\times E$ where $B$ is a generic abelian variety of dimension $g$ with polarization of type determined by $\lambda$ and $E$ is an elliptic curve ($\lambda$ also encodes the structure of the isogeny). Along each $S'_\lambda$, possibly after passing to a generically finite cover $S''_\lambda\rightarrow S'_\lambda$, we have a universal family $\mathcal{B}_{S''_\lambda}\rightarrow S''_\lambda$ of abelian varieties of dimension $g$, a morphism of families of abelian varieties
\[
p_\lambda:\mathcal{A}_{S''_\lambda}\rightarrow \mathcal{B}_{S''_\lambda}
\]
and, denoting by $\Sigma'_{\mathcal{A}_{S''_\lambda}} $ the fibered product $\Sigma'_{\mathcal{A}_{S'_\lambda}} \times_{S'_\lambda}S''_\lambda$, we have
\[
p_\lambda(\Sigma'_{\mathcal{A}_{S''_\lambda}})\subset \Sigma_{\mathcal{B}_{S''_\lambda}}
\]
by axiom~\eqref{2.1.1} of Definition~\ref{definat}.

\begin{lemm}\label{lefini}
If $\Sigma_B\subset B$ is a proper subset for a very general
abelian variety $B$ of dimension $g$, the morphism $p_{\lambda,\Sigma}:=p_{\lambda\mid \Sigma'_{\mathcal{A}_b}}:\Sigma'_{\mathcal{A}_b}\rightarrow \mathcal{B}_b$ is generically finite onto its image for any point $b$ of $S''_\lambda$.
\end{lemm}

\begin{proof}
As $p_\lambda(\Sigma'_{\mathcal{A}_b})\subset \Sigma_{\mathcal{B}_b}$ for $b\in S''_\lambda$, and since we know by assumption that $\dim \Sigma_{\mathcal{B}_b}<g$, we conclude that $\Sigma'_{\mathcal{A}_b}\subset \mathcal{A}_b$ is a proper algebraic subset for very general $b\in S''_\lambda$, hence also for very general $b\in S'$. For very general $b\in S'$, the cycle class $[\Sigma'_{\mathcal{A}_b}]\in H^{2l}(\mathcal{A}_b,\mathbb{Q})$, where $l:=\codim \Sigma'_{\mathcal{A}_b}$, is a nonzero multiple $\mu \theta^l$ of $\theta^l$ because the latter generates the space of degree-$2l$ Hodge classes of a very general abelian variety with polarizing class $\theta$ (see~\cite{mattuck}). Using the fact that $l>0$, we thus conclude that $p_{\lambda*}([\Sigma'_{\mathcal{A}_b}])=p_{\lambda*}(\mu \theta^l)$ is nonzero in $H^{2l-2}(\mathcal{B}_b,\mathbb{Q})$, and, since $\Sigma'_{\mathcal{A}_b}$ is irreducible by construction, it follows that $p_{\lambda,\Sigma}$ is generically finite on its image.
\end{proof}

Lemma~\ref{lefini} implies statement~\eqref{2.4.a} of Proposition~\ref{propourtheogenpirola} and also, applied to the case where $\Sigma_{{B}}\subset {B}$ is countable while $\dim {B}>0$, statement~\eqref{2.4.a'}. We now concentrate on statement~\eqref{2.4.b} and thus assume that $\dim B>\dim \Sigma_B>0$ for a very general abelian variety $B$ of dimension $g$. With the same notation as before, we want to show that $\dim \Sigma'_{\mathcal{A}_b}<\dim \Sigma_{\mathcal{B}_b}$ which, using as before Lemma~\ref{lefini}, means that $p_\lambda(\Sigma'_{\mathcal{A}_b})$ is not a component of $\Sigma_{\mathcal{B}_b}$ when $b\in S''_\lambda$ is general, for sufficiently general $\lambda$.

\begin{lemm}\label{lebounded}
In the situation above, the set of varieties (of dimension $k'=\dim \Sigma'_{\mathcal{A}_b}$ by Lemma~\ref{lefini})
\[
\Sigma'_{\mathcal{A}_b,p_\lambda}:=p_\lambda(\Sigma'_{\mathcal{A}_b})
\]
and morphisms $ p_{\lambda,\Sigma}: \Sigma'_{\mathcal{A}_b}\rightarrow \Sigma'_{\mathcal{A}_b,p_\lambda}$, for all $\lambda$, is bounded up to birational transformations.
\end{lemm}

The meaning of the above lemma will be made clear from its proof. One way to state boundedness here is to say that the union of the irreducible components of the subvarieties
\[
\Sigma'_{\mathcal{A}_b}\times_{\Sigma'_{\mathcal{A}_b,p_\lambda}}\Sigma'_{\mathcal{A}_b}\subset
\Sigma'_{\mathcal{A}_b}\times\Sigma'_{\mathcal{A}_b}
\]
dominating both factors $\Sigma'_{\mathcal{A}_b}$ have bounded degree for the originally given polarization on $\mathcal{A}_b$.

\begin{proof}[Proof of Lemma~\ref{lebounded}]
Recall that $\Sigma'_{\mathcal{A}_b,p_\lambda}=p_\lambda(\Sigma'_{\mathcal{A}_b})$ is contained in $ \Sigma_{\mathcal{B}_b}\subset \mathcal{B}_b$, hence is a proper subvariety of a very general abelian variety $\mathcal{B}_b$ of dimension $g$ with polarization of certain type, and $\Sigma'_{\mathcal{A}_b}\subset \mathcal{A}_b$ is the specialization of a subvariety (of codimension at least $2$ by Lemma~\ref{lefini}) of a very general abelian variety of dimension $g+1$ at points $b$ which form a Zariski dense subset of $S$. In both cases, it follows that the Gauss maps $g_A$ of $\Sigma'_{\mathcal{A}_b}\subset \mathcal{A}_b$ and $g_B$ of $\Sigma'_{\mathcal{A}_b,p_\lambda}\subset \mathcal{B}_b$, which take their respective values in $G(k',g+1)=\Grass(k',T_{\mathcal{A}_b,0_{\mathcal{A}_b}})$ and $G(k',g)=\Grass(k',T_{\mathcal{B}_b,0_{\mathcal{B}_b}})$, are generically finite on their images. We have a commutative diagram
\begin{equation}\label{numerodiag}
\begin{aligned}
\xymatrix{ \Sigma'_{\mathcal{A}_b}\ar@{-->}[r]^-{g_A}\ar@{-->}[d]^{p_{\lambda,\Sigma}} &G(k',g+1)\ar@{-->}[d]^{\pi_\lambda}\\
\Sigma'_{\mathcal{A}_b,p_\lambda}\ar@{-->}[r]^-{g_B} &G(k',g) }
\end{aligned}
\end{equation}
where all the maps are rational maps and the rational map $\pi_\lambda:G(k',g+1)\dashrightarrow G(k',g)$ is induced by the linear map $dp_\lambda: T_{\mathcal{A}_b,0_{\mathcal{A}_b}}\rightarrow T_{\mathcal{B}_b,0_{\mathcal{B}_b}}$ which is also the quotient map $T_{\mathcal{A}_b,0_{\mathcal{A}_b}}\rightarrow T_{\mathcal{A}_b,0_{\mathcal{A}_b}}/T_{E,0_E}$. We observe here that the density of the countable union of the $S_\lambda$ in $S$ has the following stronger version.

\begin{lemm}\label{ledensetEisog}
The points $[T_{E,0_E}]\in \mathbb{P}(T_{\mathcal{A}_b,0_{\mathcal{A}_b}})$, for $b\in S_\lambda$ are Zariski dense (and even dense for the usual topology) in the projectivized bundle $\mathbb{P}(T_{\mathcal{A}/S,0_{\mathcal{A}}})$ over $S$.
\end{lemm}

\begin{proof}
View a $g$-dimensional abelian variety $A$ as a complex torus $\Gamma_{\mathbb{C}}/(\Gamma_{1,0}\oplus\Gamma)$, where $\Gamma$ is a fixed lattice of rank $2g$ and $\Gamma_{\mathbb{C}}$ is the corresponding complexified vector space. The lattice $\Gamma$ is equipped with a skew nondegenerate pairing $\langle\,\cdot\,,\cdot\,\rangle$ determining the polarization and when $A$ deforms along $S$, the complex subspace $\Gamma_{1,0}\subset \Gamma_{\mathbb{C}}$, which naturally identifies with $T_{A,0_A}$ varies in an open set of the Lagrangian Grassmannian of $(\Gamma_{\mathbb{C}},\langle\,\cdot\,,\cdot\,\rangle)$. The condition that $A_b$ is isogenous to $B\times E$, that is $b\in S_\lambda$ for some $\lambda$, is equivalent to the fact that the complex line
\[
T_{E}\subset T_{A,0_A}=\Gamma_{1,0}\subset \Gamma_{\mathbb{C}}
\]
satisfies the property that the plane $ T_{E}\oplus \overline{ T_{E}}\subset \Gamma_{\mathbb{C}}$, which is defined over $\mathbb{R}$, is defined over $\mathbb{Q}$. The lemma then follows from the density for the usual topology of the set of planes defined over $\mathbb{Q}$ in the Grassmannian of real planes in $ \Gamma_{\mathbb{R}}$.
\end{proof}

This lemma shows that the projection $\pi_\lambda$ above is thus generic when $b$ is taken generic in a general $S_\lambda$ so that for general $\lambda$ and generic $b\in S_\lambda$, the composition $\pi_\lambda\circ g_A$ is generically finite as is $g_A$ and, up to shrinking $S'$ if necessary, its graph deforms in a flat way over the space of parameters (namely a Zariski open set of $\mathbb{P}(T_{\mathcal{A}/S,0_{\mathcal{A}}})$).

This is now finished because we first restrict to the Zariski dense open set $U$ of $\mathbb{P}(T_{\mathcal{A}_{S/B},0_{\mathcal{A}}})$ where the rational map $\pi_\lambda\circ g_A$ is generically finite and its graph deforms in a flat way, and then there are finitely many generically finite covers of $U$ parameterizing a factorization of the rational map $\pi_\lambda\circ g_A$. Since the diagram~\eqref{numerodiag} shows that there is a factorization of $\pi_\lambda\circ g_A$ as
\[
\Sigma'_{\mathcal{A}_b}\stackrel{p_{\lambda,\Sigma}}{\rightarrow }\Sigma'_{\mathcal{A}_b,p_\lambda}\stackrel{g_B}{\rightarrow }G(k',g),
\]
we conclude that all the maps $\Sigma'_{\mathcal{A}_b}\stackrel{p_{\lambda,\Sigma}}{\rightarrow }\Sigma'_{\mathcal{A}_b,p_\lambda}$ are, up to birational equivalence of the target, members of finitely many families of generically finite dominant rational maps $\psi:\mathcal{A}_b\dashrightarrow Y_b$.
\end{proof}

As a corollary, we conclude using the density of the union of the sets $S'_\lambda$ that there is, up to replacing $S'$ by a generically finite cover, a family of $k'$-dimensional varieties $\Sigma''_{\mathcal{A}_{S'}}$, together with a dominant generically finite rational map
\begin{equation}\label{eqpfam}
p: \Sigma'_{\mathcal{A}_{S'}}\dashrightarrow \Sigma''_{\mathcal{A}_{S'}}
\end{equation}
identifying birationally to $p_{\lambda,\Sigma}:\Sigma'_{\mathcal{A}_{b}}\dashrightarrow
\Sigma'_{\mathcal{A}_{b,p_\lambda}}$ along each $S'_\lambda$.

We now finish the proof by contradiction. Assume that $k'=k$. Then $\Sigma'_{\mathcal{A}_b,p_\lambda}$ must be a component $\Gamma_b$ of $\Sigma_{\mathcal{B}_b}$. In particular it does not depend on the elliptic curve $E$. Restricting to a dense Zariski open set $S''$ of $S'$ is necessary, we can assume that we have desingularizations
\begin{equation}\label{eqpfam1}
\widetilde{\Sigma'_{\mathcal{A}_{S''}}},\quad \widetilde{\Sigma''_{\mathcal{A}_{S''}}}
\end{equation}
with smooth fibers over $S''$, together with a surjective generically finite morphism $\tilde{p}:\widetilde{\Sigma'_{\mathcal{A}_{S''}}}\rightarrow \widetilde{\Sigma''_{\mathcal{A}_{S''}}}$ over $S'$. Let $\tilde{j}:\Sigma'_{\mathcal{A}_b}\rightarrow \mathcal{A}_b$ be the natural map and consider the morphism of abelian varieties
\[
\tilde{p}_*\circ \tilde{j}^*:\Pic^0(\mathcal{A}_b)\rightarrow \Pic^0(\widetilde{\Sigma''_{\mathcal{A}_{b}}})
\]
which is defined at a general point of $S''$. This morphism is nonzero because when $b\in S''_\lambda$ for some $\lambda$, it is injective modulo torsion on $\Pic^0(\mathcal{B}_b)$ (which maps by the pull-back $p_\lambda^*$ to $ \Pic^0(\mathcal{A}_b)$ with finite kernel). Indeed, by the projection formula, denoting by $\tilde{j'}:\widetilde{\Sigma''_{\mathcal{A}_{b}}}\rightarrow B$ the natural map, we have the equality of maps from $\Pic^0(\mathcal{B}_b)$ to $\Pic^0(\widetilde{\Sigma''_{\mathcal{A}_{b}}})$:
\[
(\tilde{p}_*\circ \tilde{j}^*)_{\mid \Pic^0(\mathcal{B}_b)}=
((\tilde{p}_{\lambda,\Sigma})_*\circ \tilde{j}^*)_{\mid \Pic^0(\mathcal{B}_b)}=
((\tilde{p}_{\lambda,\Sigma})_*\circ(\tilde{p}_{\lambda,\Sigma})^*\circ \tilde{j'}^*=\deg p_{\lambda,\Sigma} \,\tilde{j'}^*.
\]
(We note here that the morphism $\tilde{j'}^*:\Pic^0(\mathcal{B}_b)\rightarrow\Pic^0(\widetilde{\Sigma''_{\mathcal{A}_{b}}})$ has finite kernel because $\dim \im \tilde{j'}=k>0$. It is at this point that we use the assumption that $\dim \Sigma_B>0$.)

As the abelian variety $\Pic^0(\mathcal{A}_b)$ is simple at a very general point $b$ of $S''$, the nonzero morphism $(\tilde{p}_{\lambda,\Sigma})_*\circ \tilde{j}^*$ must be injective up to torsion. But then, by specializing at a point $b$ of $S''_\lambda$, where $\lambda$ is chosen in such a way that $S''_\lambda=S''\cap S'_\lambda$ is non-empty, we find that this morphism is injective up to torsion on the component $\Pic^0(E_b)$ of $\Pic^0(\mathcal{A}_b)$. We can now fix the abelian variety $\mathcal{B}_b$ and deform the elliptic curve $E_b$. We then get a contradiction, because we know that the variety $\widetilde{\Sigma''_{\mathcal{A}_{b}}}$ depends (at least birationally) only on $\mathcal{B}_b$ and not on $E_b$, so that its Picard variety cannot contain a variable elliptic curve $E_b$.
\end{proof}


\section{Proof of Theorems~\ref{theochowring} and~\ref{coroorbit}\tMK{(\ref{1.4.4})}}\label{secproofcoriv}

\subsection{Dimension estimate}\label{secdimestimate}

Recall that for an abelian variety $A$ and a nonnegative integer $k$, we denote by $A_k\subset A$ the set of points $x\in A$ such that $(\{x\}-\{0_A\})^{*k}=0$ in $\CH_0(A)$. Let us first for completeness prove Lemma~\ref{lefourier}, which says that the isogeny between $A$ and $\widehat{A}$ provided by a polarization induces a surjection with finite fibers between $A_k$ and the subset $\widehat{A}_k$ of divisors homologous to $0$ in $A$ such that $D^k=0$ in $\CH^k(A)$.

\begin{proof}[Proof of Lemma~\ref{lefourier}]
We use Beauville's formulas in~\cite[Proposition~6]{beauvillefourier}. We get in particular the following equality:
\begin{equation}\label{eqbeau}
\frac{\theta^{g-k}}{(g-k)!}D_x^k=\frac{\theta^g}{g!}*\gamma(x)^{*k},
\end{equation}
where
\[
\gamma(x):=\{0_A\}-\{x\}+\frac{1}{2}(\{0_A\}-\{x\})^{*2}+\dots+\frac{1}{g}(\{0_A\}-\{x\})^{*g}=-\log(\{x\}) \text{ in }\CH_0(A).
\]
Here the logarithm is taken with respect to the Pontryagin product $*$ and the expansion is finite because $0$-cycles of degree~$0$ are nilpotent for the Pontryagin product. If $(\{0_A\}-\{x\})^{*k}=0$, then $\gamma(x)^{*k}=0$ and thus ${\theta^{g-k}}D_x^k=0$ in $\CH_0(A)$ by~\eqref{eqbeau}. This implies that $D_x^k=0$ by the following lemma.

\begin{lemm}\label{lesurintA}
Let $(A,\theta)$ be a polarized abelian variety of dimension $g$, and $D\in \Pic^0(A)\otimes\mathbb{Q}=\CH^1(A)_{\hom}$ be a divisor homologous to $0$, where $k\leq g$. Then $\theta^{g-k}D^k=0$ in $\CH_0(A)$ if and only if $D^k=0$ in $\CH^k(A)$.
\end{lemm}

\begin{proof}[Sketch of proof]
This lemma is very similar to~\cite[Proposition~4.6]{bloch} and could be in fact also reduced to it by a Fourier transform argument. We choose a smooth ample curve $C\subset A$ which is a complete intersection of smooth hypersurfaces of class proportional to $\theta$, inducing by the sum map a surjective morphism
\[
\sigma: C^g\rightarrow A.
\]
We then compute $\sigma^*(D^k) \text{ in }\CH(C^g)$. Using the fact that for a divisor $D$ homologous to $0$ on $A$, one has $\sigma^*D=\sum_{i=1}^g\pr_i^*(D_{\mid C})$, one concludes that $\sigma^*(D^k)=0 \text{ in }\CH^k(C^g)$ if and only if $\sigma_k^*(D^k)=0$ in $\CH^k(C^k)$, where $\sigma_k:C^k\rightarrow A$ is the sum map for $C^k$. One then computes $\sigma^*(\theta^{g-k}D^k)\in\CH_0(C^g)$ and shows that it vanishes if and only if $\sigma_k^*(D^k)=0\in \CH^k(C^k)$. One uses for this the fact that $\sigma^*\theta$ is a divisor of the form $\sum_{i=1}^{g}\pr_i^*H+\lambda\sum_{i\not=j}\Delta_{ij}$, for some divisor $H$ on $C$, where $\Delta_{ij}\subset C^g$ is the diagonal $\{c_i=c_j\}$ of $C^g$.
\end{proof}

Conversely, if $D_x^k=0$, then $\gamma(x)^{*k}=0$ by~\eqref{eqbeau}. But then also $(\{0_A\}-\{x\})^{*k}=0$ because $\{x\}=\exp(-\gamma(x))$. (Again $\exp(-\gamma(x))$ is a polynomial in $\gamma(x)$, hence well defined since $\gamma(x)$ is nilpotent for the $*$-product, see~\cite{bloch}.)
\end{proof}

The following proves item~\eqref{1.8.1} of Theorem~\ref{theochowring}.

\begin{prop}\label{proestimatechowring}
For $k>0$, one has $\dim  A_k\leq k-1$.
\end{prop}

\begin{proof}
Let $g:=\dim A$ and let $\Gamma^{\Pont}_k$ be the codimension $g$ cycle of
$A\times A$ such that
\[
(\Gamma^{\Pont}_{k })_*(x)=(\{x\}-\{0_A\})^{*k}
\]
for any $x\in A$. As $(\{x\}-\{0_A\})^{*k}=\sum_{i=0}^k(-1)^{k-i}\binom{k}{i}\{ix\}$, we can take
\begin{equation}\label{eqstarkx}
\Gamma^{\Pont}_k=\sum_{i=0}^k(-1)^{k-i}\binom{k}{i}\Gamma_{i},
\end{equation}
where $\Gamma_i\subset A\times A$ is the graph of the map $m_i$ of multiplication by $i$. Let us compute $(\Gamma_k^{\Pont})^*\eta$ for any holomorphic form on $A$.

\begin{lemm}\label{lekform}
One has $(\Gamma^{\Pont}_k)^*\eta=0$ for any holomorphic form $\eta$ of degree $<k$ on $A$, and $(\Gamma_k^{\Pont})^*\eta=k!\eta$ for any holomorphic form of degree $k$ on $A$.
\end{lemm}

\begin{proof}
Indeed $m_i^*\eta=i^d\eta$, where $d=\deg \eta$.
By~\eqref{eqstarkx}, the lemma is thus equivalent to
\begin{enumerate}
\item \label{p3.3.1} $\sum_{i=0}^k(-1)^{k-i}\binom{k}{i} i^d=0,\,d<k$,
\item \label{p3.3.2} $\sum_{i=0}^k(-1)^{k-i}\binom{k}{i} i^k=k!$.
\end{enumerate}

From the formula $\sum_{i=0}^k(-1)^{k-i}\binom{k}{i}X^i=(X-1)^k$, we get that the $d$-th derivative of the polynomial $\sum_{i=0}^k(-1)^{k-i}\binom{k}{i}X^i$ at $1$ is $0$ for $d<k$, and is equal to $k!$ for $d=k$. This implies~\eqref{p3.3.1} by induction on $d$, using the fact that $i(i-1)\dots(i-d+1)-i^d$ is a polynomial of degree $d-1$ in $i$, and then~\eqref{p3.3.2} by the same argument.
\end{proof}

This lemma implies Proposition~\ref{proestimatechowring}. Indeed, by Mumford's theorem~\cite{mumford}, one has $(\Gamma^{\Pont}_k)^*\eta_{\mid W}=0$ for any smooth algebraic variety $W$ contained in $A_k$ and any holomorphic form $\eta$ of positive degree on $A$, and in particular for any holomorphic $k$-form on $A$. By Lemma~\ref{lekform}, we conclude that, for any $W$ as above and for any holomorphic form $\eta$ of degree $k$ on $A$, we have $\eta_{\mid W}=0$. Applying this to the regular locus $W$ of any component of $A_k$, one concludes that $\dim A_k<k$.
\end{proof}

\subsection{Proof of Theorem~\ref{theochowring}}\label{secnatdef}

The following result is almost obvious.

\begin{lemm}\label{leAknat}
For any integer $k\geq 0$, the family of subsets $A_k\subset A$ defined in~\eqref{eqAk} is naturally defined in the sense of Definition~\ref{definat}.
\end{lemm}

\begin{proof}
It is known that the set $A_k\subset A$ is a countable union of closed algebraic subsets.
Using the fact that for a morphism $f:A\rightarrow B$ of abelian varieties,
\[
f_*:\CH_0(A)\rightarrow \CH_0(B)
\]
is compatible with the Pontryagin product, we conclude that $f(A_k)\subset B_k$. Finally, given a family $\pi:\mathcal{A}\rightarrow S$ of abelian varieties, the set of points $x\in \mathcal{A}$ such that $(\{x\}-\{0_{\mathcal{A}_b}\})^{*k}=0$ in $\CH_0(\mathcal{A}_b)$, where $b=\pi(x)$, is a countable union of closed algebraic subsets of $\mathcal{A}$ whose fiber over $b\in S$ coincides set-theoretically with $\mathcal{A}_{b,k}$.
\end{proof}

\begin{proof}[Proof of Theorem~\ref{theochowring}]
The theorem follows from Proposition~\ref{proestimatechowring}, Lemma~\ref{leAknat}, and Theorem~\ref{thegenpirola}.
\end{proof}

\subsection{Proof of Theorem~\ref{coroorbit}\tMK(\ref{1.4.4})}\label{sectheoorbitoA}

We first prove the following proposition (cf. Proposition~\ref{theorelation}). The elegant proof given below is due to A.~Beauville and simplifies our original proof.

\begin{prop}\label{theorelation2}
Let $A$ be an abelian variety and let $x_1,\dots,\,x_k\in A$ such that
\begin{equation}\label{eqhypo}
\sum_i\{x_i\}=k\{0_A\}\,\,\, in \,\,\,\CH_0(A).
\end{equation}
Then
\begin{equation}\label{eqpoutheorrelation}
(\{x_i\}-\{0_A\})^{*k}=0 \text{ in }\CH_0(A)
\end{equation}
for all $i=1,\dots,\,k$.
\end{prop}

\begin{proof}
Letting act the isogeny of multiplication by $l$ on $A$, we get from relation~\eqref{eqhypo}
\begin{equation}\label{eqhypol}
\sum_i\{lx_i\}=k\{0_A\}\text{ in }\CH_0(A).
\end{equation}
Let $y_i:=\{x_i\}-\{0_A\}\in\CH_0(A)$. The relations~\eqref{eqhypol} and the binomial formulas show inductively that for any $l>0$
\begin{equation}\label{eqhypoly}
\sum_iy_i^{*l}=0 \text{ in }\CH_0(A).
\end{equation}
The Newton formulas expressing any polynomial symmetric function of degree $\leq m$ in the $y_i$ as a polynomial in the power sums $\sum_iy_i^{*l}$ for $l\leq m$ then imply that the elementary symmetric functions $s_l(y_1,\dots,\,y_k)$ vanish for $l=1,\dots,\,k$. On the other hand, the $y_i$ are roots of the polynomial (where the product is Pontryagin product in $\CH_0(A)$)
\[
P(t):=\prod_{i=1}^k (t-y_i).
\]
The vanishing of the symmetric functions $s_l(y_1,\dots,\,y_k)$ for $l=1,\dots,\,k$ then gives $y_i^{*k}=0$, that is, \eqref{eqpoutheorrelation}.
\end{proof}

Proposition~\ref{theorelation2} says that if $\{x_1\}+\dots+\{x_k\}=k\{0_A\}$ in $\CH_0(A)$, then $x_i\in A_k$. The locus swept out by the orbit $|k\{0_A\}|$ is thus contained in $A_k$. We thus deduce from Theorem~\ref{theochowring} the following corollary:

\begin{coro}[cf. Theorem~\ref{coroorbit}$\MK$\eqref{1.4.4}]
For any abelian variety $A$, the locus swept out by the orbit $|k\{0_A\}|$ has dimension $\leq k-1$. For a very general abelian variety $A$ of dimension $g\geq 2k-1$, the orbit $|k\{0_A\}|$ is countable.
\end{coro}

In this statement, the locus swept out by the orbit $|k\{0_A\}|$ is the set of points $x\in A$ such that there exists a cycle $x+Z'$, with $Z'$ effective of degree $k-1$, which belongs to $|k\{0_A\}|$. The dimension of this locus can be much smaller than the dimension of the orbit itself, as shown by the examples of orbits contained in subvarieties $C^{(k)}\subset A^{(k)}$ for some curve $C$.


\section{Proof of Theorem~\ref{coroorbit}\tMK(\ref{1.4.1})}\label{secdimorbit}

We give in this section the proof of item~\eqref{1.4.1} in Theorem~\ref{coroorbit}. We first recall the statement:

\begin{theo}\label{ledimorbit}
Let $A$ be an abelian variety. The dimension of any orbit
$|Z|\subset A^{(k)}$ for rational equivalence is at most $k-1$.
\end{theo}

\begin{proof}
We will rather work with the inverse image $\widetilde{|Z|}$ of the orbit $|Z|$ in $A^k$. By Mumford's theorem~\cite{mumford}, for any holomorphic $i$-form $\alpha$ on $A$ with $i>0$, one has, along the regular locus $\widetilde{|Z|}_{{\reg}}$ of $\widetilde{|Z|}$:
\begin{equation}\label{eqmumford}
\sum_{j=1}^k\pr_j^*\alpha|_{ \widetilde{|Z|}_{\reg}}=0,
\end{equation}
where the $\pr_j:A^k\rightarrow A$ are the various projections. Let $x=(x_1,\dots,x_k)\in \widetilde{|Z|}_{\reg}$ and let $V:=T_{\widetilde{|Z|}_{\reg},x}\subset W^k$, where $W=T_{A,x}=T_{A,0_A}$. One has $\dim V=\dim |Z|$ and~\eqref{eqmumford} says that:
\begin{enumerate}
\renewcommand{\theenumi}{$\ast$}
\item \label{*} for any $\alpha\in \bigwedge^iW^*$ with $i>0$, one has $(\sum_j\pr_j^*\alpha)_{\mid V}=0$.
\end{enumerate}
Theorem~\ref{ledimorbit} thus follows from the following proposition.
\end{proof}

\begin{prop}\label{ledimorbitinfinitesimal}
Let $W$ be a vector space, and let $V\subset W^k$ be a vector subspace
satisfying property~\eqref{*}. Then $\dim V\leq k-1$.
\end{prop}

\begin{rema}
If $\dim W=1$, the result is obvious, as $V\subset W^k_0\subset W$, where,
denoting by $\sigma$ the sum map, $W^k_0:=\Ker (\sigma: W^k\rightarrow W)$. If $\dim W=2$, the result follows from the fact that, choosing a generator $\eta$ of $\bigwedge^2W^*$, the $2$-form $\sum_j\pr_j^*\eta$ is nondegenerate on $W^k_0$ (which has dimension~$2k-2$). A subspace $V$ satisfying~\eqref{*} is contained in $W^k_0$ and totally isotropic for this $2$-form, hence has dimension $r\leq k-1$.
\end{rema}

\begin{proof}[Proof of Proposition~\ref{ledimorbitinfinitesimal}]
The group $\Aut W$ acts on $W^k$, with induced action on $\Grass(r,W^k)$ preserving the set of $r$-dimensional vector subspaces $V\subset W^k$ satisfying condition~\eqref{*}. Choose a $\mathbb{C}^*$-action on ${W}$ with finitely many fixed points $e_1,\dots, e_n$, where $n=\dim W$. The fixed points $[V]\in \Grass(r,W^k)$ under the induced action of $\mathbb{C}^*$ on the Grassmannian are of the form $V=\langle A_1e_1,\dots, A_n e_n\rangle$, where $A_i\subset (\mathbb{C}^k)^*$ are vector subspaces, with $r=\sum_i\dim A_i$. It suffices to prove the inequality $r\leq k-1$ at such a fixed point, which we do now. The spaces $A_i$ have to satisfy the following conditions:
\begin{enumerate}
\renewcommand{\theenumi}{$\ast\ast$}
\item \label{**}
For any nonempty subset $I=\{i_1,\dots,\,i_s\}\subset \{1,\dots,n\}$ and for any choices of $\lambda_l\in A_{i_l}$, $l=1,\dots,\,s$,
\[
\sum_{j=1}^k(\lambda_1\dots\lambda_s)(f_j)=0,
\]
where $\{f_1,\dots, \,f_k\}$ is the natural basis of $\mathbb{C}^k$.
\end{enumerate}

A better way to phrase condition~\eqref{**} is to use the (standard) pairing $\langle\,\cdot\,,\cdot\,\rangle$ on $(\mathbb{C}^k)^*$ given by
\[
\langle \alpha,\beta\rangle=\sum_{j=1}^k\alpha(f_j)\beta(f_j).
\]
Condition~\eqref{**} when there are only two nonzero spaces $A_i$ is the following
\begin{align}\label{eqdeux1}
\forall \alpha\in A_1,\,\beta\in A_2,\,\,\,\,\,\, \langle\alpha,\beta\rangle=0\\\label{equn}
\langle \alpha,e\rangle=0,\,\,\langle e,\beta\rangle=0,
\end{align}
where $e$ is the vector $(1,\dots,1)\in (\mathbb{C}^k)^*$. Indeed, the case $s=2$ in~\eqref{**} provides~\eqref{eqdeux1} and the case $s=1$ in~\eqref{**} provides~\eqref{equn}. The fact that the pairing $\langle\,\cdot\,,\cdot\,\rangle$ is nondegenerate on $(\mathbb{C}^k)^*_0:=e^{\perp}$ immediately implies that $\sum_i\dim A_i\leq k-1$ when only two of the spaces $A_i$ are nonzero. By the above arguments, the proof of Proposition~\ref{ledimorbitinfinitesimal} is finished using the following lemma:

\begin{lemm}\label{lepourproduitdeformes}
Let $A_1,\dots,\,A_n$ be vector subspaces of
$(\mathbb{C}^k)^*$ satisfying conditions~\eqref{**}. Then $\sum_i\dim A_i\leq k-1$.
\end{lemm}

\begin{proof}
We will use the following result:

\begin{lemm}\label{ledeuxA}
Let $A$ and $B$ be two vector subspaces
of $\mathbb{C}^k$ satisfying the following conditions:
\begin{enumerate}
\item \label{4.5.1} For any $a=(a_i)\in A$ and $b=(b_i)\in B$, one has $\sum_ia_ib_i=0$.
\item \label{4.5.2} For any $a=(a_i)\in A$ and $b=(b_i)\in B$, one has $\sum_ia_i=0$ and $\sum_ib_i=0$.
\end{enumerate}
Then $\dim (A\cdot B +A+B)\geq \dim A+\dim B$, where $A\cdot B$ is the vector subspace of $\mathbb{C}^k$ generated by the elements $(a_ib_i)$, $a=(a_i)\in A,\,b=(b_i)\in B$.
\end{lemm}

Let us first show how Lemma~\ref{ledeuxA} implies Lemma~\ref{lepourproduitdeformes}. Indeed, we can argue inductively on the number $n$ of spaces $A_i$. As already noticed, Lemma~\ref{lepourproduitdeformes} is easy when $n=2$. Assuming the statement is proved for $n-1$, let $A_1,\dots,\,A_n$ be as in Lemma~\ref{lepourproduitdeformes} and let $A'_1=A_1,\,\dots,\,A'_{n-2}=A_{n-2}$ and $A'_{n-1}=A_{n-1}\cdot A_n+A_{n-1}+A_n$. Then the set of spaces $A'_1,\dots,\,A'_{n-1}$ satisfies conditions~\eqref{**}, and on the other hand Lemma~\ref{ledeuxA} applies to the pair $(A,B)=(A_{n-1},A_n)$ as they satisfy the desired conditions by~\eqref{**}. Hence we have $\dim A'_{n-1}\geq \dim A_{n-1}+ \dim A_{n}$ and by induction on $n$, we obtain $\sum_{i=1}^{n-2}\dim A'_{i}+\dim A'_{n-1}\leq k-1$. Hence $\sum_{i=1}^{n}\dim A_{i}\leq k-1$.
\end{proof}

\begin{proof}[Proof of Lemma~\ref{ledeuxA}]
Recalling that $e=(1,\dots,1)\in \mathbb{C}^k$,
consider the affine subspaces $A_1:= e+A$ and $B_1:=e+B$ of $ \mathbb{C}^k$. Under the conditions~\eqref{4.5.1} and~\eqref{4.5.2}, the componentwise multiplication map
\begin{align*}
\mu:\;\; A_1\times B_1&\rightarrow \mathbb{C}^k\\
((a_i),(b_i))&\mapsto(a_ib_i)
\end{align*}
has image in the affine space $\mathbb{C}^k_1:=e+\mathbb{C}^k_0$, where $\mathbb{C}^k_0=e^{\perp}$, and more precisely it generates the affine space $e+A+B+A\cdot B\subset e+\mathbb{C}^k_0$. It thus suffices to show that the dimension of the algebraic set $\im \mu$ is at least $\dim A+\dim B$. Lemma~\ref{ledeuxA} is thus implied by the following:

\begin{enonce}{Claim}\label{claim}
The map $\mu$ has finite fibers near the point $(e,e)\in A_1\times B_1$.
\end{enonce}

The proof of the claim is as follows: Suppose $\mu$ has a positive-dimensional fiber passing through $(e,e)$. We choose an irreducible curve contained in the fiber, passing through $(e,e)$ and with normalization $C$. The curve $C$ admits rational functions $\sigma_1,\dots,\,\sigma_k$ mapping it to $A_1$ such that the functions $\frac{1}{\sigma_i}$ map $C$ to $B_1$.
The conditions~\eqref{4.5.1} and~\eqref{4.5.2} say that
\[
\sum_i\sigma'_i(s)\frac{1}{\sigma_i(t)}=0
\]
as a function of $(s,t)$ for any choice of points $x,\,y\in C$ and local coordinates $s,\,t$ near $x$, resp. $y$, on $C$. We now take $x=y$ and choose for $x$ a pole (or a zero) of one of the $\sigma_l$. We assume that the local coordinate $s$ is centered at $x$ and write
$\sigma_i(s)=s^{d_i}f_i(s)$, where $f_i$ is a holomorphic function of $s$ which is nonzero at $0$. We then get
\begin{equation}\label{eqpourfonction}
\sigma'_i(s)\frac{1}{\sigma_i(t)}=d_i \frac{s^{d_i-1}}{t^{d_i}} \phi_i(s,t)+\frac{s^{d_i}}{t^{d_i}}\psi_i(s,t),
\end{equation}
where $\phi_i(s,t)$ is holomorphic in $s,\,t$ and takes value $1$ at $(0,0)$ and $\psi_i(s,t)$ is holomorphic in $s,\,t$. Restricting to a curve $D\subset C\times C$ defined by the equation $s=t^l$ for some chosen $l\geq2$, the function $(\sigma'_i(s)\frac{1}{\sigma_i(t)})_{\mid D}$ has order $l(d_i-1)-d_i=(l-1)d_i-l$ and first nonzero coefficient in its Laurent expansion equal to $d_i$. These orders are different for distinct $d_i$ and the vanishing $\sum_i\sigma'_i(s)\frac{1}{\sigma_i}(t)=0$ is then clearly impossible: indeed, by pole order considerations, for the minimal negative value $d$ of $d_i$, hence minimal value of the numbers $(l-1)d_i-l$, the first nonzero coefficient in the Laurent expansion of $(\sigma'_i(s)\frac{1}{\sigma_i(t)})_{\mid D}$ should be also $0$. However, it is the same as for the sum $\sum_{i,\,d_i=d}(\sigma'_i(s)\frac{1}{\sigma_i}(t))_{\mid D}$, which is equal to $M_d d$, where $M_d$ is the cardinality of the set $\{i,\,d_i=d\}$, hence it is nonzero.

The claim is proved.
\end{proof}

The proof of Proposition~\ref{ledimorbitinfinitesimal} is thus finished.
\end{proof}

\subsection{An alternative proof of Theorem~\ref{coroorbit}\tMK(\ref{1.4.4})}

As a first application, let us give a second proof of Theorem~\ref{coroorbit}$\MK$\eqref{1.4.4}. The general dimension estimate of Theorem~\ref{coroorbit}$\MK$\eqref{1.4.1} implies that the locus swept out by the orbit of $|k0_A|$ is of dimension $\leq k-1$ for any abelian variety $A$. This locus is clearly naturally defined. Hence by Theorem~\ref{thegenpirola}$\MK$\eqref{2.3.2}, it is countable for a very general abelian variety of dimension $\geq 2k-1$.


\section{Proof of Theorem~\ref{coroorbit}\tMK(\ref{1.4.2}) and~\ref{coroorbit}\tMK(\ref{1.4.3})}

We will prove the following result by induction on $l\in\{0,\dots,k\}$.

\begin{prop}\label{proinduction}
For $g\geq 2^{l}(2k-1)+(2^{l}-1)(k-2)$ and for $A$ a very general abelian variety
of dimension $g$, any $0$-cycle of the form $(k-l)\{0_A\}+Z$, with $Z\in A^{(l)}$, has countable orbit.
\end{prop}

The case $l=0$, which is the first step in our inductive proof, is precisely Theorem~\ref{coroorbit}$\MK$\eqref{1.4.4} which was proved in Section~\ref{sectheoorbitoA}.

The case $l=k$ implies Theorem~\ref{coroorbit}$\MK$\eqref{1.4.2}, because any effective zero-cycle $Z$ of degree $k$ can be translated using one of its points so as to pass through $0_A$. The translation by a given element acts on $A$, hence on $A^{(k)}$, preserving the dimension of the orbit. Hence there exists a zero-cycle of degree $k$ with positive-dimensional orbit for rational equivalence in $A$ if and only if there exists a zero-cycle of degree $k$ of the form $0_A+Z'$, with $Z'$ effective of degree $k-1$, with a positive-dimensional orbit for rational equivalence in $A$.

The case $l=k-2$ is Theorem~\ref{coroorbit}$\MK$\eqref{1.4.3}.

It thus only remains to prove Proposition~\ref{proinduction}. For clarity, let us write up the detail of the first induction step: let
\begin{equation}\label{eqnomsigma1}
\Sigma_1(A)\subset A
\end{equation}
be the set of points $x\in A$ such that the orbit $|(k-1)\{0_A\}+\{x\}|\subset A^{(k)}$ is positive-dimensional. The set $\Sigma_1(A)$ is a countable union of closed algebraic subsets of $A$. We would like to show that $\Sigma_1(A)$ is naturally defined in the sense of Definition~\ref{definat}, but there is a small difficulty here: suppose that $p:A\rightarrow B$ is a morphism of abelian varieties and let $|Z|\subset A^{(k)} $ be a positive-dimensional orbit for rational equivalence on $A$. Then $p_*(|Z|)\subset B^{(k)} $ could be zero-dimensional. In the case where $Z=(k-1)\{0_A\}+\{x\}$, this a priori prevents us from proving that $\Sigma_1(A)$ satisfies axiom~\eqref{2.3.2} of Definition~\ref{definat}.

This problem can be circumvented using the following lemma whose main contents in fact already appeared in the course of the proof of Theorem~\ref{thegenpirola} (see Lemma~\ref{lefini}). Let $\mathcal{A}\rightarrow S$ be a family of abelian varieties of dimension $g$ which is locally complete at the generic point. This means that we fixed a polarization type $\lambda$ and the moduli map $S\rightarrow\mathcal{A}_{g,\lambda}$ is dominant.

\begin{lemm}\label{generic}
Let $\mathcal{W}\subset \mathcal{A}$ be a closed algebraic subset which is
flat over $S$ of relative dimension $k'$. Then:
\begin{enumerate}
\item \label{5.2.1} For any $b\in S$, any morphism $p: \mathcal{A}_b\rightarrow B$ of abelian varieties with $\dim B\geq k'$, $p(\mathcal{W}_b)\subset B$ has dimension $k'$.
\item \label{5.2.2} Assume $k'>0$. For any $b\in S$, any morphism $p: \mathcal{A}_b\rightarrow B$ of abelian varieties with $\dim B>0$, $p(\mathcal{W}_b)\subset B$ has positive dimension.
\end{enumerate}
\end{lemm}

\begin{proof}
Let's first prove Statement~\eqref{5.2.1}. Indeed, the locally constant class $[\mathcal{W}_b]\in H^{2g-2k'}(\mathcal{A}_b,\mathbb{Q})$ must be a nonzero multiple of $\theta_\lambda^{g-k'}$, since for very generic $b\in S$, these are the only nonzero Hodge classes on $\mathcal{A}_b$ by~\cite{mattuck}. Using our assumption that $\dim B\geq k'$, we thus get that $p_*([\mathcal{W}_b])\not=0$ in $H^{2\dim B-2k'}(B,\mathbb{Q})$, which implies that $\dim p(\mathcal{W}_b)=k'$.

Statement~\eqref{5.2.2} is obtained as an application of~\eqref{5.2.1} in the case $k'=1$. By replacing $\mathcal{W}_b$ by a complete intersections in $\mathcal{W}_b$, one can assume that $\dim \mathcal{W}_b=1$. Then one has $\dim \mathcal{W}_b\leq\dim B$ if $\dim B>0$, so~\eqref{5.2.1} applies.
\end{proof}

In the following corollary, the orbits for rational equivalence of $0$-cycles of $X$ are taken in $X^l$ rather than $X^{(l)}$. Given a family $\mathcal{A}\rightarrow S$, we denote by $\mathcal{A}^{l/S}$ its $l$-th power over $S$. A family of positive-dimensional orbits for rational equivalence in the fibers is a closed algebraic subset $\mathcal{W}\subset \mathcal{A}^{l/S}$ for some $l$, which is of positive relative dimension over $S$, and whose fibers $\mathcal{W}_b$ for $b\in S$ are orbits for rational equivalence in the fibers $\mathcal{A}_b$.

\begin{coro}\label{coropourtheoorbitgen}
Let $\mathcal{W}\subset \mathcal{A}^{l/S}$ be a family of positive-dimensional orbits for rational equivalence in the fibers. Then, up to shrinking $S$ if necessary, for any $b\in S$, any morphism $p: \mathcal{A}_b\rightarrow B$ of abelian varieties, where $B$ is an abelian variety of dimension $>0$, $p^l(\mathcal{W}_b)\subset B^l$ is a positive-dimensional orbit of $B$.
\end{coro}

Here $p^l:\mathcal{A}_b^l\rightarrow B^l$ is the morphism induced by $p$.

\begin{proof}
Indeed, by specialization, $\mathcal{W}_b $ is a positive-dimensional
orbit for rational equivalence in $\mathcal{A}_b^{l}$. Up to shrinking $S$, we can assume that the restrictions $\pi_{\mid \pr_i(\mathcal{W})}: \pr_i(\mathcal{W})\rightarrow S$ are flat for all $i$. Our assumption is that for one $i$, $\pr_i(\mathcal{W})$ has positive relative dimension over $S$. Lemma~\ref{generic}$\MK$\eqref{5.2.2} then implies that $\pr_i(p^l(\mathcal{W}_b))$ has positive dimension, so that $p^l(\mathcal{W}_b)$ is a positive-dimensional orbit for rational equivalence of $0$-cycles of $B$.
\end{proof}

\begin{proof}[Proof of Proposition~\ref{proinduction}]
Let $A$ be a very general abelian variety. This means that for some generically complete family $\pi: \mathcal{A}\rightarrow S$ of polarized abelian varieties, $A$ is isomorphic to the fiber of $\pi$ over a very general point of $S$. As $A$ is very general, the locus $\Sigma_1(A)$ introduced in~\eqref{eqnomsigma1} is the specialization of the corresponding relative locus $\Sigma_1(\mathcal{A}/S)$ of $\mathcal{A}$, and more precisely, of the union of its components dominating $S$. For any fiber $\mathcal{A}_b$, let us define the deformable locus $\Sigma_1(A)_{\rmdef}$ as the one which is obtained by specializing to $\mathcal{A}_b$ the union of the components of the relative locus $\Sigma_1(\mathcal{A}/S)$ dominating $S$. For a very general abelian variety $A$, $\Sigma_1(A)=\Sigma_1(A)_{\rmdef}$ by definition.

Corollary~\ref{coropourtheoorbitgen} essentially says that this locus is naturally defined. This is not quite true because the definition of $\Sigma_1(A)_{\rmdef}$ depends on choosing a family $\mathcal{A}$ of deformations of $A$ (that is, a polarization on $A$). In the axioms of Definition~\ref{definat}, we thus should work, not with abelian varieties but with polarized abelian varieties. Axiom~\eqref{2.1.1} should be replaced by its family version, where $\mathcal{A}\rightarrow S$ is locally complete, $S'\subset S$ is a subvariety, $f:\mathcal{A}_{S'}\rightarrow \mathcal{B}$ is a morphism of abelian varieties over $S'$, and $\mathcal{B}\rightarrow {S'}$ is locally complete. We leave to the reader the task of proving that Theorem~\ref{thegenpirola} extends to this context.

Assume now $g\geq 2k-1$. Then $\Sigma_1(A)$, hence a fortiori $\Sigma_1(A)_{\rmdef}$, is different from $A$. Indeed, otherwise, for any $x\in A$, $(k-1)\{0_A\}+\{x\}$ has positive-dimensional orbit, hence taking $x=0_A$, we get that $k\{0_A\}$ has positive-dimensional orbit, contradicting Theorem~\ref{coroorbit}$\MK$\eqref{1.4.4}. Theorem~\ref{thegenpirola}$\MK$\eqref{2.3.1} then implies that for $g\geq 2(2k-1)-1$, the set $\Sigma_1(A)_{\rmdef}$ is countable. Hence there are only countably many positive-dimensional orbits of the form $|(k-1)\{0_A\}+\{x\}|$ and the locus they sweep-out forms by Corollary~\ref{coropourtheoorbitgen} a naturally defined locus in $A$, which is of dimension $\leq k-1$ by Theorem~\ref{ledimorbit}. It follows by applying Theorem~\ref{thegenpirola}$\MK$\eqref{2.3.3} that for $g\geq 2(2k-1)+k-2$, this locus itself is countable, that is, all the orbits $|(k-1)\{0_A\}+\{x\}|$ are countable for $A$ very general.

The general induction step works exactly in the same way, introducing the locus $\Sigma_l(A)\subset A$ of points $x_l\in A$ such that $(k-l)0_A+x_1+\dots+x_l$ has a positive-dimensional orbit for rational equivalence in $A$ for some points $x_1,\dots,\,x_{l-1}\in A$.
\end{proof}

\section{Further discussion}\label{secdiscussion}

It would be nice to improve the estimates in our main theorems. As already mentioned in the introduction, none of them seem to be optimal. Let us introduce a naturally defined locus (or the deformation variant of that notion used in the last section) whose study should lead to a proof of Conjecture~\ref{conj}.

\begin{defi}
The locus $Z_A\subset A$ of positive-dimensional normalized orbits of degree $k$
is the set of points $x\in A$ such that for some degree-$k$ zero-cycle $Z=x+Z'$, with $Z'$ effective, one has
\[
\dim |Z|>0,\,\,\sigma(Z)=0.
\]
\end{defi}

Here $\sigma: A^{(k)}\rightarrow A$ is the sum map. It is constant along orbits under rational equivalence. This locus (or its ``def'' version discussed in the previous proof) is naturally defined. Note also that by definition it is either of positive dimension or empty. The main remaining question is to estimate the dimension of this locus, at least for very general abelian varieties. Conjecture~\ref{conj} would follow from:

\begin{conj}\label{conjchow}
If $A$ is a very general abelian variety, the locus $Z_A\subset A$ of positive-dimensional normalized orbits of degree $k$ has dimension $\leq k-1$.
\end{conj}

Conjecture~\ref{conjchow} is true for $k=2$. Indeed, in this case the normalization condition reads $Z=\{x\}+\{-x\}$ for some $x\in A$. The positive-dimensional normalized orbits of degree~$2$ are thus also positive-dimensional orbits of points in the Kummer variety $K(A)=A/\pm \Id$ of $A$. These orbits are rigid because on a surface in $K(A)$ swept out by a continuous family of such orbits, any holomorphic $2$-form on $K(A)$ should vanish while $\Omega^2_{K(A)_{\reg}}$ is generated by its sections.

It would be tempting to try to estimate the dimension of the locus of positive-dimensional normalized orbits of degree $k$ for any abelian variety. Unfortunately, the following example shows that this locus can be the whole of $A$.

\begin{exam}\label{exanom}
Let $A$ be an abelian variety which has a degree-$k-1$
positive-dimensional orbit $Z\subset A^{(k-1)}$. Then for each $x\in A$,
\[
x+Z:=\{\{x_1+x\}+\dots+\{x_{k-1}+x\},\,\,\,\,\{x_1\}+\dots+\{x_{k-1}\}\in Z\}
\]
is also a positive-dimensional orbit. The element $\sum_ix_i$ of $A$ is constant along $Z$
and for fixed $x$, $Z$, the set
\[
\{\{x_1+x\}+\dots+\{x_{k-1}+x\}+\{-\sum_ix_i-(k-1)x\}\}\subset A^{(k)}
\]
is a positive-dimensional \emph{normalized} orbit of degree $k$.
In this case, the locus of positive-dimensional normalized orbits of degree $k$ of $A$ is the whole of $A$.
\end{exam}

Nevertheless, we can observe the following small evidence for Conjecture~\ref{conjchow}.

\begin{lemm}\label{ledimorbitavechypot}
Let $O\subset A^k$ be a closed irreducible algebraic subset which is covered by components of positive-dimensional normalized orbits of degree $k$. Let $Z\in O_{\reg}$ and assume
the positive-dimensional orbit $O_Z$ passing through $Z$ has a tangent vector $(u_1,\dots, u_k)$ such that the vector space $\langle u_1,\dots,u_k\rangle\subset T_{A,0_A}$ is of dimension $k-1$. Then the locus of $A$ swept out by $O$, namely the union of the sets $pr_i(O)\subset A$ has dimension $\leq k-1$.
\end{lemm}

Note that $k-1$ is the maximal possible dimension of the vector space $\langle u_1,\dots,u_k\rangle$ because $\sum_iu_i=0$. The example~\ref{exanom} is a case where the vector space $\langle u_1,\dots,u_k\rangle$ has dimension~$1$.

Applying Theorem~\ref{thegenpirola}$\MK$\eqref{2.3.2}, Conjecture~\ref{conjchow} in fact implies the following.

\begin{conj}\label{conjchowdeux}
If $A$ is a very general abelian variety of dimension $\geq 2k-1$, $A$ has no positive-dimensional orbits of degree $k$.
\end{conj}

This is a generalization of Conjecture~\ref{conj}, because a $k$-gonal curve ${j}:{C}\rightarrow A,\,D\in W_k^1(C)$ provides a positive-dimensional orbit ${j}_*|D|$ of degree $k$ of $A$.

We discussed in this paper only the applications to gonality. The case of higher dimensional linear systems would be also interesting to investigate. In a similar but different vein, the following problem is intriguing:

\begin{ques}\label{questionplancurve}
Let $A$ be a very general abelian variety. Is it true that there are no a smooth plane curves $C$ admitting a nonconstant morphism to $ A$ ? Equivalently, is it true that there is no surjective morphism $J(C)\rightarrow A$ of abelian varieties, with $C$ a smooth plane curve?
\end{ques}

If the answer to the above question is affirmative, one could get examples of surfaces of general type which are not birational to a normal surface in $\mathbb{P}^3$. Indeed, take a surface whose Albanese variety is a general abelian variety as above. If $S$ is birational to a normal surface $S'$ in $\mathbb{P}^3$, there are plenty of smooth plane curves in $S'$, which clearly map nontrivially to $\Alb S$, which would be a contradiction.

